% !PRINTED 429
\weberpage{429}

We wish to prove that primitive roots exist for every prime $n$, and to
determine their number; and in doing so we follow the same path as with
the roots of unity.

Let $a$ and $b$ be two divisors of $n-1$ which are relatively prime
to each other, and let $\alpha$ be an $a$th, $\beta$ a $b$th
primitive congruence root of $n$. If we set $ab=c$, then
$\gamma=\alpha\beta$ is a primitive $c$th congruence root, and
conversely every primitive $c$th congruence root can also be represented
in the form $\alpha\beta$.

For, first: $(\alpha\beta)^c = \alpha^c\beta^c \equiv 1 \pmod{n}$;
and secondly: if $(\alpha\beta)^h = \alpha^h\beta^h \equiv 1 \pmod{n}$,
then also $\alpha^{hb}\beta^{hb}\equiv 1$, hence $\alpha^{hb}\equiv 1$,
hence $hb$ is divisible by $a$, hence $h$ too is divisible by $a$; and in
the same way one concludes that $h$ must be divisible by $b$, and hence
also by $ab=c$; that is, $\alpha\beta$ is a primitive $c$th
congruence root.

If conversely $\gamma$ is a primitive $c$th congruence root, then we
determine the positive integers $x, y$ so that
\weq{10}{a y + b x \equiv 1 \pmod{c}}
holds, by first determining $x$ from $bx \equiv 1 \pmod{a}$ and then $y$ from
$y \equiv \dfrac{1-bx}{a} \pmod{b}$. Then $y$ is relatively prime
to $b$ and $x$ relatively prime to $a$. Accordingly
\[
\gamma \equiv \gamma^{bx}\gamma^{ay} = \alpha\beta \pmod{n}
\]
and $\gamma^{bx} = \alpha$ is a primitive $a$th
congruence root, $\gamma^{ay} = \beta$ a primitive $b$th
congruence root (because $\alpha^h = \gamma^{bxh}$ can be $\equiv 1$
only when $h$ is divisible by $a$).

It remains to be shown that the products $\alpha\beta$ are all distinct
from one another, that is, that the congruence
\weq{11}{\alpha\beta \equiv \alpha'\beta' \pmod{n},}
when $\alpha, \alpha'$; $\beta, \beta'$ are primitive $a$th and
$b$th congruence roots, can be satisfied only when
\[
\alpha \equiv \alpha', \quad \beta \equiv \beta' \pmod{n}
\]
holds. But this follows by (10), if one raises (11) to the power
\[
bx \equiv 1 \pmod{a}
\]
from which $\alpha \equiv \alpha'$ results, and then also
$\beta \equiv \beta'$.

If, then, we denote the number of primitive $a$th
congruence roots by $\varphi(a)$, it follows from what has been proved that
\weq{12}{\varphi(c) = \varphi(a)\varphi(b).}
